\section{Joint refinement and its bijective game}
\label{game:section}

We first fix the refinement convention and prove the game characterization
used in the reduction.  The standard joint convention and parameter shift are described in
\citep[Section~2.2 and Theorem~1]{KieferNeuen2019}. We prove the precise
histogram and empty-position statement, including the order in which
a bijection is announced.

\begin{definition}[Joint-update Weisfeiler--Leman refinement]
\label{def:joint-wl}
Let $X$ be a finite simple graph and let $k\geq 2$.  For
$\mathbf v=(v_1,\ldots,v_k)\in V(X)^k$, write
$\mathbf v[i\leftarrow z]$ for the tuple obtained by replacing its $i$th
coordinate by $z$.  The initial color $\chi_0^X(\mathbf v)$ records, for every
ordered pair of coordinates, whether their entries are equal and whether
they are adjacent.  Recursively, set
\begin{equation}
\label{game:joint-update}
 \chi_{s+1}^X(\mathbf v)=
 \left(\chi_s^X(\mathbf v),
 \left\{\!\!\left\{
 \bigl(\chi_s^X(\mathbf v[1\leftarrow z]),\ldots,
       \chi_s^X(\mathbf v[k\leftarrow z])\bigr):z\in V(X)
 \right\}\!\!\right\}\right).
\end{equation}
The braces in the second component denote a multiset.  The recursive names
of colors are common to all graphs under comparison.  We write
$G\equiv_k H$ if, at every round $s\geq 0$, the multisets of colors on
$V(G)^k$ and $V(H)^k$ agree.
\end{definition}

In particular, all $k$ coordinate replacements associated with one vertex
$z$ remain grouped together in~\eqref{game:joint-update}.  This is the joint
update, rather than an update using $k$ separate multisets.

\begin{lemma}[Capping the dimension]
\label{lem:dimension-cap}
Let $G,H$ have the same positive order $n$.  For every $k\geq n$ with
$k\geq 2$, the relation $G\equiv_k H$ holds if and only if $G$ and $H$ are
isomorphic.  Consequently, when $n\geq 2$ and $k\geq 2$,
\[
 G\equiv_k H \quad\Longleftrightarrow\quad
 G\equiv_{\min\{k,n\}}H.
\]
\end{lemma}

\begin{proof}
An isomorphism acts coordinatewise on tuples and preserves their initial
colors.  Induction in~\eqref{game:joint-update} shows that it preserves every
later color as well, and hence preserves each color histogram.

Conversely, suppose $k\geq n$ and the initial histograms agree.  Choose a
tuple $\mathbf v\in V(G)^k$ containing every vertex of $G$.  There is a tuple
$\mathbf w\in V(H)^k$ with the same initial color.  The common equality
pattern shows that the rule $v_i\mapsto w_i$ is a well-defined bijection
between the $n$ distinct entries of $\mathbf v$ and the $n$ distinct entries
of $\mathbf w$.  These are all vertices of the two graphs.  The adjacency
part of the initial color shows that this bijection is an isomorphism.
Applying the same assertion with $k=n$ proves the displayed equivalence
when $k\geq n$; when $k<n$, it is an identity.
\end{proof}

\begin{proposition}[Uniform upper bound]
\label{prop:membership}
Both \(\WL\) and \(\SubWL\) belong to \(\EXPTIME\).
More precisely, after checking their respective input conditions,
an input of common graph order \(n\ge2\) and total bit length \(L\)
can be decided in time
$2^{O(n\log n)}L^{O(1)}$, with absolute constants in the exponents.
\end{proposition}

\begin{proof}
First check that the input consists of the prescribed two adjacency matrices
and a binary integer $k\geq 2$.  Check matrix lengths before allocating data
for any claimed order.  Symmetry, zero diagonal, and entries in \(\{0,1\}\) are checked in
polynomial time. For \(\SubWL\), also check positive equal orders,
maximum degree three, and connectedness. Reject a failed condition.
For \(\WL\), unequal orders give different tuple counts and are
rejected; at equal orders zero or one the simple graphs are necessarily
isomorphic and are accepted. At order one the promised subcubic graphs
are also isomorphic.

For $n\geq 2$, compare the binary integer $k$ with $n$ and replace it by
$\min\{k,n\}$, using \Cref{lem:dimension-cap}.  This comparison takes
polynomial time in $L$ even when the value of the supplied integer is
enormous.  In the rest of the algorithm, therefore, $2\leq k\leq n$.

Maintain tables for $V(G)^k$ and $V(H)^k$, and give their union a common
system of integer color labels.  There are
\[
 A=2n^k
\]
table entries.  To perform a round, form the replacement-color vector for
each of the $n$ possible replacement vertices at each entry, sort these
vectors with multiplicities, and include the old label in the resulting
signature.  Sorting all signatures assigns the next common labels.  The
replacements in a graph's table use only vertices of that graph; mixed
tuples are neither needed nor introduced.

Every new partition refines the old one because the old label is part of
the signature.  A strict refinement increases the number of classes, which
never exceeds $A$.  If a round makes no split, its labels are a bijective
renaming of the old classes, and substituting those names into the next
signatures cannot create a split.  Thus at most $A$ rounds, including a
possible final no-split round, establish stabilization.  Compare the two
histograms initially and after every round, rejecting as soon as they
differ.  If they agree at stabilization, every later update merely renames
the same common classes, so all later histograms also agree.

An integer label uses $O(1+\log A)$ bits.  A replacement vector has $k$
such labels, and a signature has $n$ such vectors together with an old
label.  Explicit table access, sorting, comparison, and counting can be
implemented in time bounded by a fixed polynomial in $A+n+k+L$ per round,
even with sequential access.  Since $k\leq n$ and $n\geq 2$, multiplying
by at most $A$ rounds gives
\[
 (A+n+k+L)^{O(1)}=2^{O(n\log n)}L^{O(1)}.
\]
The graph order is polynomially bounded by the length of its explicit
adjacency matrix.  This is a uniform exponential-time bound in $L$.
\end{proof}

We now define the game directly, so that no convention about the timing of
a pebble move is implicit.  For equal-order graphs $G,H$ and an integer
$q\geq 1$, the game has $q$ named \emph{pair slots}.  An occupied slot holds
one vertex of $G$ and one vertex of $H$.  Initially all slots are vacant.
Spoiler may discard any occupied slots.  To place a new pair, Spoiler
requests a vacant slot; Duplicator then announces a bijection
$f:V(G)\to V(H)$; finally Spoiler chooses $v\in V(G)$, and the requested
slot receives $(v,f(v))$.  The same vertex may occur in several slots.
After every move the marked pairs must form a partial isomorphism: equality
and adjacency between any two occupied slots must agree in the two graphs.
Duplicator wins if a strategy, allowed to depend on the entire finite
history, maintains this condition for every finite play.  No time bound is
imposed on either player's strategy.

\begin{lemma}[Joint refinement and announced bijections]
\label{lem:wl-game}
For graphs $G,H$ of the same positive order and every $k\geq 2$,
$G\equiv_k H$ if and only if Duplicator wins the game with $k+1$ pair slots.
\end{lemma}

\begin{proof}
We first record a symmetry of the colors.  If two tuples have equal
round-$s$ colors, applying the same permutation to their coordinates gives
tuples of equal round-$s$ colors.  At round zero this follows from the
definition of the atomic data.  For the induction step, equal next-round
colors give equal old colors and a matching of replacement vertices whose
entire old-color vectors agree.  By the induction hypothesis, applying the
coordinate permutation to each replaced tuple preserves these equalities.
Reordering the vector components by the same permutation then gives the
replacement vectors for the permuted tuples.  Their multisets agree, as
required.  Notice that the assertion concerns equality of colors, not that
an individual color name is fixed by every permutation.

Suppose first that $G\equiv_k H$.  On the union of their tuple tables, let
$\sim$ be equality in a stable common refinement partition.  For
$\mathbf v\sim\mathbf w$, stability gives a bijection
$f:V(G)\to V(H)$ such that, simultaneously for every $z\in V(G)$ and every
$i\in\{1,\ldots,k\}$,
\begin{equation}
\label{game:stable-bijection}
 \mathbf v[i\leftarrow z]\sim\mathbf w[i\leftarrow f(z)].
\end{equation}
Indeed, the stable-color vectors in the two multisets can be paired with
their multiplicities, and this pairing matches the replacement vertices
bijectively.

The stable histograms agree and are nonempty, so some related tuple pair
exists.  Append any vertex and its image under a bijection
from~\eqref{game:stable-bijection}.  The resulting pair of $(k+1)$-tuples
has the following property: deleting any one coordinate produces related
$k$-tuples, in the common order of the remaining coordinates.  Deleting
the new entry recovers the original related pair.  Deleting an old entry
gives the corresponding replacement pair up to a common permutation, to
which the preceding symmetry applies.

Duplicator keeps such a full pair of $(k+1)$-tuples as auxiliary data,
with the entries of occupied slots equal to their actual marked entries
and the entries of vacant slots regarded as virtual.  Discarding a slot
does not change this data.  When slot $j$ is requested, delete its virtual
entries to obtain a related pair of $k$-tuples, and announce the bijection
in~\eqref{game:stable-bijection} for that pair.  Overwrite slot $j$ with
Spoiler's chosen pair.  The projection omitting $j$ is unchanged.  A
projection omitting any other slot is, up to a common coordinate
permutation, one of the replacements in~\eqref{game:stable-bijection}.
Thus the auxiliary invariant persists.  Any two occupied coordinates occur
together in some $k$-projection because $k\geq 2$.  Their equality and
adjacency therefore agree, proving that the strategy wins.

Conversely, fix a winning strategy for Duplicator.  We prove by induction
on $s$ that at every reached history with exactly $k$ occupied slots, the
two tuples in any common ordering of those slots have equal round-$s$
colors.  For $s=0$ this is the partial-isomorphism condition.  Suppose the
assertion holds at round $s$, and consider such a history with ordered
tuples $\mathbf v,\mathbf w$.  Request the one vacant slot and let $f$ be
the announced bijection.  For each choice of $z\in V(G)$, place
$(z,f(z))$.  For each $i$, there is a continuation that discards the old
$i$th slot and orders the remaining $k$ slots by putting the new slot in
its position.  The induction hypothesis, which applies to every reached
history, gives
\[
 \chi_s^G(\mathbf v[i\leftarrow z])
 =\chi_s^H(\mathbf w[i\leftarrow f(z)]).
\]
The bijection $f$ was announced before $z$ was chosen and is the same for
all of these alternatives.  It matches the entire replacement vectors,
and the old tuple colors also agree by induction.  Equation
\eqref{game:joint-update} now gives equality at round $s+1$.

Finally, fill $k$ fixed slots from empty in a fixed order, with no
discards, following the strategy.  This procedure defines a map
$F:V(G)^k\to V(H)^k$.  It is bijective: given a target tuple, invert the
first announced bijection to recover the first source entry; the recovered
history determines the second announced bijection, which recovers the
second entry; continue in this way.  This inversion remains valid when
the strategy depends on its full history.  For every $s$, the induction
just proved says that $F$ preserves round-$s$ colors.  Hence the two
histograms agree at every round, proving the required equivalence.
\end{proof}

We will repeatedly use one elementary consequence of winning.  At a
reached history, a property of specified marked pairs can be established
by showing that its failure would admit a losing continuation after all
other pairs are discarded.  Such an argument establishes the property
in the original history; it does not require executing the test during a
different intended continuation, nor restoring any discarded answers.
